\documentclass[11pt,a4paper]{article}
\usepackage[utf8]{inputenc}
\usepackage[T1]{fontenc}
\usepackage[french]{babel}
\usepackage{lmodern}
\usepackage{amsmath,amssymb}
\usepackage{graphicx}

\usepackage[margin=2.5cm]{geometry}
\usepackage{fancyhdr}
\pagestyle{fancy}
\fancyhf{}
\fancyhead[L]{ETE7-TSIG2}
\fancyhead[R]{Cirou -- Païtard}
\fancyfoot[C]{\thepage}
\renewcommand{\headrulewidth}{0.4pt}

\begin{document}

\begin{titlepage}
  \centering
  {\large ENSEIRB-MATMECA}\\[0.2cm]
  {\large Filière Télécommunications -- 2\textsuperscript{e} année}\\[0.2cm]
  {\large ETE7-TSIG2 : Filtres numériques, estimation et applications}\\[3cm]

  \rule{\linewidth}{0.5pt}\\[0.4cm]
  {\LARGE\bfseries Rapport Projet Signaux}\\[0.3cm]
  {\large Extraction de marqueurs en vue d'une classification de signaux}\\
  \rule{\linewidth}{0.5pt}\\[3cm]

  {\large Marco \textsc{Cirou} \quad -- \quad Maxence \textsc{Païtard}}\\[1cm]
  {\large Encadrant : Éric \textsc{Grivel}}

  \vfill
  {\large Octobre 2026}
\end{titlepage}

\tableofcontents
\newpage

\section{Introduction}

\section{Signaux Aléatoires}
\subsection{Processus résultant du filtrage d'un bruit blanc gaussien}
\subsubsection{1ère réalisation}
Soit un bruit blanc gaussien $u(n)$, centré et de variance $\sigma_u^2$, filtré par un filtre dont les quatre zéros de module $\rho = 0{,}95$ sont placés aux fréquences normalisées : $\nu_1 = -4/9$, $\nu_2 = -1/4$, $\nu_3 = 1/4$ et $\nu_4 = 4/9$. Les zéros sont donc de la forme : $z_k = \rho\, e^{j2\pi\nu_k}$ \\
On a donc : \\
\begin{equation}
  B(z) = \prod_{k=1}^{4} \left( 1 - z_k\, z^{-1} \right)
       = \sum_{i=0}^{4} b_i z^{-i},
  \qquad
  x(n) = \sum_{i=0}^{4} b_i\, u(n-i),
  \label{eq:ma4}
\end{equation} \\
On remarque que la moyenne de ce processus est nulle : $\mathbb{E}[x(n)] = 0$. Par ailleurs, son autocorrélation ne dépend que de $k$ :
\begin{equation}
    \forall |k| > 4,
    \qquad
    r_x(k) = \sigma_u^2 \sum_{i=0}^{4-|k|} b_i b_{i+|k|}.
    \label{eq:r_x}
\end{equation}
Ainsi il s'agit d'un processus stationnaire au sens large. Le filtre est un RIF car sa fonction de transfert
$B(z)$ ne possède que des zéros et la sortie ne dépend que d'un nombre fini
d'échantillons d'entrée, la fonction d'autocorrélation
est nulle pour tout $\forall |k| > 4$.\\
Le processus $x(n)$ est donc un processus à moyenne mobile d'ordre 4. \\ \\
La dénsité spectrale de puissance s'écrit : \\
\begin{equation}
  S_x(\nu) = \sigma_u^2 \left| B\!\left(e^{j2\pi\nu}\right) \right|^2.
  \label{eq:dsp_ma}
\end{equation}
Les zéros étant proches du cercle unité, $S_x(\nu)$ présente deux creux marqués en
$\nu_1$ et $\nu_2$, visibles sur la figure~\ref{fig:ma}.

\begin{figure}[h]
  \centering
  \includegraphics[width=\linewidth]{ma.pdf}
  \caption{Processus MA(4) : réalisation de $N = 2000$ échantillons (gauche) et densité
  spectrale de puissance théorique en dB (droite).}
  \label{fig:ma}
\end{figure}

\subsubsection{2\textsuperscript{de} réalisation}

De manière analogue, le bruit blanc $u(n)$ est désormais filtré par un filtre dont les
quatre pôles, de module $\rho = 0{,}99$, s'écrivent $p_k = \rho\, e^{j2\pi\nu_k}$, avec
$\nu_1 = -3/7$, $\nu_2 = -1/3$, $\nu_3 = 1/3$ et $\nu_4 = 3/7$. On a donc :
\begin{equation}
  H(z) = \frac{1}{\prod_{k=1}^{4} \left( 1 - p_k\, z^{-1} \right)}
       = \frac{1}{1 + \sum_{i=1}^{4} a_i z^{-i}},
  \qquad
  x(n) = u(n) - \sum_{i=1}^{4} a_i\, x(n-i).
  \label{eq:ar4}
\end{equation}
L'échantillon $x(n)$ dépend désormais de ses quatre échantillons passés : il s'agit d'un
processus autorégressif d'ordre 4.

Contrairement au cas précédent, il s'agit ici d'un filtre RII. Sa
stabilité n'est donc plus garantie : elle est assurée ici car tous les pôles sont à
l'intérieur du cercle unité ($\rho < 1$). Le processus est alors stationnaire au sens
large, une fois le régime transitoire dû aux conditions initiales écoulé. Les pôles étant
proches du cercle unité, ce régime est plus long que dans le cas MA. Les premiers
échantillons de la réalisation sont donc écartés. La moyenne reste nulle, mais
l'autocorrélation n'est plus nulle au-delà d'un certain rang : elle vérifie les équations
de Yule-Walker, $r_x(k) = -\sum_{i=1}^{4} a_i\, r_x(k-i)$ pour $k \geq 1$.

Par analogie avec \eqref{eq:dsp_ma}, la densité spectrale de puissance s'écrit :
\begin{equation}
  S_x(\nu) = \frac{\sigma_u^2}{\left| A\!\left(e^{j2\pi\nu}\right) \right|^2},
  \qquad A(z) = 1 + \sum_{i=1}^{4} a_i z^{-i}.
  \label{eq:dsp_ar}
\end{equation}
Les creux du cas MA laissent place à des pics de résonance en $\nu_3 = 1/3$ et
$\nu_4 = 3/7$, visibles sur la figure~\ref{fig:ar}.

\begin{figure}[h]
  \centering
  \includegraphics[width=\linewidth]{ar.pdf}
  \caption{Processus AR(4) : réalisation de $N = 2000$ échantillons (gauche) et densité
  spectrale de puissance théorique en dB (droite).}
  \label{fig:ar}
\end{figure}

\subsubsection{Représentations temporelles}

\subsubsection{Périodogrammes}
Daniell : \\ \\
Le périodogramme de Daniell fait la valeur absolue de la transformée de Fourrier au carré du signal sur une fenêtre, et se déplace d'un terme à chaque occurence.\\
Ainsi, on a : \\
\begin{equation}
    \hat{S}_D(\nu_k) = \frac{1}{2m+1} \sum_{i=-m}^{m} I_N(\nu_{k+i}).
  \label{eq:daniell}
\end{equation}

Bartlett : \\ \\

\begin{equation}
    \hat{S}_x^{(k)}(w) = \frac{1}{M} abs(\sum_{n=0}^{M-1} x_k[n]e^{-jwn}).
  \label{eq:daniell}
\end{equation}

Welch \\ \\
bh
Corrélogramme \\ \\


\section{Méthode à mettre en oeuvre}

\subsection{Étape n°1}

\subsection{Étape n°2}

\subsection{Étape n°3}

\subsection{Étape n°4}

\section{Applications}


\section{Organisation des Séances}

\section{Annexes}

L'ensemble des codes Matlab est fourni dans l'archive
\texttt{Cirou\_Paitard\_TSIG2.zip}, jointe au rapport. Le script principal
\texttt{main.m} permet de reproduire l'ensemble des figures présentées. Chaque fonction
précise en commentaire d'en-tête ses entrées, ses sorties et son objet.
\end{document}
